% GENERATED FILE. Edit canonical lessons, then run npm run generate:books.
% canonical-source-hash: fnv1a64:db475c9a8b86639a
% canonical-lessons: KA-C100-map, KA-C100-ticket, KA-C100-bus, KA-C100-station, KA-C100-near

\chapter{Help Me Find It}
\label{ch:ka-help-me-find-it}
\hlchaptermodality{100}

\begin{chapteropening}
\textbf{By the end of this chapter:} \emph{I can ask for a map, a ticket, a bus and a station, and say it is near.}

Complete the last lesson of chapter 100: i can ask for a map, a ticket, a bus and a station, and say it is near.
\end{chapteropening}

\section[nakṣe]{\kn{ನಕ್ಷೆ} (nakṣe) — a map}

\label{lesson:KA-C100-map}



\begin{quote}
Before the new one: say the Kannada for a stomach ache, then the Kannada for weight.
\end{quote}

\subsection*{You'll want to know: \kn{ನಕ್ಷೆ}}
\textbf{\kn{ನಕ್ಷೆ}} (\emph{nakṣe}) — \textquotedblleft{}a map\textquotedblright{}. The \textbf{\kn{ಕ್ಷ}} is \emph{k} and \emph{ṣ} joined, the same pair as in \textbf{\kn{ಕ್ಷಮಿಸಿ}}.

\textbf{\kn{ನಕ್ಷೆ} \kn{ತೋರಿಸಿ}} (\emph{nakṣe tōrisi}) — \textquotedblleft{}please show me the map\textquotedblright{}.

\begin{grammarlens}[title={What you've built}]
Something to point at when words run out.
\end{grammarlens}

\subsection*{Your turn}
\begin{itemize}
\raggedright
  \item \emph{Say it:} \emph{nakṣe}
  \item \emph{Say it:} \emph{nakṣe}, once more
  \item \emph{Say it:} say \emph{nakṣe tōrisi}
  \item \emph{Recall:} say the Kannada for a stomach ache, then the Kannada for weight, then say \emph{nakṣe} again
\end{itemize}

\begin{tcolorbox}[breakable,colback=teal!4,colframe=teal!35!black,title={Before you move on}]
What does \kn{ನಕ್ಷೆ} mean? (\textquotedblleft{}A map\textquotedblright{}.) Say it once more.
\end{tcolorbox}
\section[ṭikeṭ]{\kn{ಟಿಕೆಟ್} (ṭikeṭ) — a ticket}

\label{lesson:KA-C100-ticket}



\begin{quote}
Before the new one: say the Kannada for weight, then the Kannada for a map.
\end{quote}

\subsection*{You'll want to know: \kn{ಟಿಕೆಟ್}}
\textbf{\kn{ಟಿಕೆಟ್}} (\emph{ṭikeṭ}) — \textquotedblleft{}a ticket\textquotedblright{}, borrowed from English. Kannada hears the English \emph{t} as the curled \textbf{\kn{ಟ}}.

\textbf{\kn{ಒಂದು} \kn{ಟಿಕೆಟ್} \kn{ಕೊಡಿ}} (\emph{ondu ṭikeṭ koḍi}) — \textquotedblleft{}one ticket, please\textquotedblright{}.

\begin{grammarlens}[title={What you've built}]
An English word, heard with Kannada ears.
\end{grammarlens}

\subsection*{Your turn}
\begin{itemize}
\raggedright
  \item \emph{Say it:} \emph{ṭikeṭ}
  \item \emph{Say it:} \emph{ṭikeṭ}, once more
  \item \emph{Say it:} say \emph{ondu ṭikeṭ koḍi}
  \item \emph{Recall:} say the Kannada for weight, then the Kannada for a map, then say \emph{ṭikeṭ} again
\end{itemize}

\begin{tcolorbox}[breakable,colback=teal!4,colframe=teal!35!black,title={Before you move on}]
What does \kn{ಟಿಕೆಟ್} mean? (\textquotedblleft{}A ticket\textquotedblright{}.) Say it once more.
\end{tcolorbox}
\section[bas]{\kn{ಬಸ್} (bas) — a bus}

\label{lesson:KA-C100-bus}



\begin{quote}
Before the new one: say the Kannada for a map, then the Kannada for a ticket.
\end{quote}

\subsection*{You'll want to know: \kn{ಬಸ್}}
\textbf{\kn{ಬಸ್}} (\emph{bas}) — \textquotedblleft{}a bus\textquotedblright{}. The \textbf{‑\kn{್}} at the end keeps the last \emph{s} bare, with no vowel after it.

\begin{grammarlens}[title={What you've built}]
The everyday way to travel.
\end{grammarlens}

\subsection*{Your turn}
\begin{itemize}
\raggedright
  \item \emph{Say it:} \emph{bas}
  \item \emph{Say it:} \emph{bas}, once more
  \item \emph{Say it:} say \emph{bas}
  \item \emph{Recall:} say the Kannada for a map, then the Kannada for a ticket, then say \emph{bas} again
\end{itemize}

\begin{tcolorbox}[breakable,colback=teal!4,colframe=teal!35!black,title={Before you move on}]
What does \kn{ಬಸ್} mean? (\textquotedblleft{}A bus\textquotedblright{}.) Say it once more.
\end{tcolorbox}
\section[nildāṇa]{\kn{ನಿಲ್ದಾಣ} (nildāṇa) — a station, a stop}

\label{lesson:KA-C100-station}



\begin{quote}
Before the new one: say the Kannada for a ticket, then the Kannada for a bus.
\end{quote}

\subsection*{You'll want to know: \kn{ನಿಲ್ದಾಣ}}
\textbf{\kn{ನಿಲ್ದಾಣ}} (\emph{nildāṇa}) — \textquotedblleft{}a station, a stop\textquotedblright{}, from \textbf{\kn{ನಿಲ್ಲು}} (\emph{nillu}), \textquotedblleft{}to stop, to stand\textquotedblright{}, which you know.

\textbf{\kn{ಬಸ್} \kn{ನಿಲ್ದಾಣ} \kn{ಎಲ್ಲಿ}?} (\emph{bas nildāṇa elli?}) — \textquotedblleft{}where is the bus stop?\textquotedblright{}

\begin{grammarlens}[title={What you've built}]
Where the bus stands.
\end{grammarlens}

\subsection*{Your turn}
\begin{itemize}
\raggedright
  \item \emph{Say it:} \emph{nildāṇa}
  \item \emph{Say it:} \emph{nildāṇa}, once more
  \item \emph{Say it:} ask \emph{bas nildāṇa elli?}
  \item \emph{Recall:} say the Kannada for a ticket, then the Kannada for a bus, then say \emph{nildāṇa} again
\end{itemize}

\begin{tcolorbox}[breakable,colback=teal!4,colframe=teal!35!black,title={Before you move on}]
What does \kn{ನಿಲ್ದಾಣ} mean? (\textquotedblleft{}A station, a stop\textquotedblright{}.) Say it once more.
\end{tcolorbox}
\section[hattira]{\kn{ಹತ್ತಿರ} (hattira) — near}

\label{lesson:KA-C100-near}



\begin{quote}
Before the new one: say the Kannada for a bus, then the Kannada for a station.
\end{quote}

\subsection*{You'll want to know: \kn{ಹತ್ತಿರ}}
\textbf{\kn{ಹತ್ತಿರ}} (\emph{hattira}) — \textquotedblleft{}near\textquotedblright{}.

\textbf{\kn{ಹತ್ತಿರ} \kn{ಇದೆಯಾ}?} (\emph{hattira ideyā?}) — \textquotedblleft{}is it near?\textquotedblright{}

\begin{grammarlens}[title={What you've built}]
The answer you hope for.
\end{grammarlens}

\subsection*{Your turn}
\begin{itemize}
\raggedright
  \item \emph{Say it:} \emph{hattira}
  \item \emph{Say it:} \emph{hattira}, once more
  \item \emph{Say it:} ask \emph{hattira ideyā?}
  \item \emph{Recall:} say the Kannada for a bus, then the Kannada for a station, then say \emph{hattira} again
\end{itemize}

\begin{tcolorbox}[breakable,colback=teal!4,colframe=teal!35!black,title={Before you move on}]
What does \kn{ಹತ್ತಿರ} mean? (\textquotedblleft{}Near\textquotedblright{}.) Say it once more.
\end{tcolorbox}
